\section*{Chapter 12 --- The hardware, the transport and the factory}
\addcontentsline{toc}{section}{Chapter 12 --- The hardware and the factory}

\solution{ex:12:1}{Watch the optimiser work}
The exercise asks for code; the outcomes are described here.
\ans{a} The loop disappears, or becomes an unconditional jump to itself. The compiler reads
\code{flag} once, sees that nothing in the loop changes it, and concludes that the condition is
constant.
\ans{b} The read moves back into the loop: one load from memory per iteration.
\ans{c} The status flag is changed by the SPI peripheral, not by the program. Without
\code{volatile} the compiler has the same right to read it once and assume that it stands still
--- and then \code{transfer()} either waits forever or does not wait at all.

\solution{ex:12:2}{Predict, change, observe}
The exercise asks for code changes; the symptoms are described here.
\ans{a} Nothing happens on the connected wires. The peripheral sits on the default mux, that is, a
different port, and the program runs on without reporting anything. The most confusing symptom
there is: no error code, no clock pulses, no clue.
\ans{b} The master leaves master mode on the first \code{select()}, since it interprets the low
\code{SS} level --- which your own code has just created --- as another master having taken over
the bus. The transfer stops immediately.
\ans{c} The clock ends up above the protocol's ceiling of 1~MHz. It may very well work on a short
cable, and that is the dangerous part: the FPGA slave guarantees nothing above it, so the faults
come sporadically and get worse with longer wires.
\ans{d} The status flag is never cleared. The next \code{transfer()} sees a flag that is still set,
returns immediately, and reads a byte that has not had time to arrive --- so all the data is
\textbf{one byte behind} through the whole transaction. It \emph{looks} like it works: bytes
arrive, the count is right, only the content is shifted, and it is easy to suspect the bit order
instead.
\ans{e} The clock pin is not driven, so the slave sees no edges and nothing is transferred. Unlike
a), \code{MOSI} and \code{SS} are active, which means a logic analyser shows activity on three of
the four wires --- and points straight at it.

\solution{ex:12:3}{The architecture test}
The exercise asks you to compile; the questions are answered here.
\ans{a} It should work. If it does not, the architecture is broken.
\ans{b} The compiler error names the file, and that file is by definition above
\code{driver/can/interface.h} and has learnt something about the hardware --- typically an
\code{#include} of \code{spi.h} or a reference to \code{ByteTransport}.
\ans{c} Exactly the same thing as before, against an array in memory: it echoes injected frames
and increments its counter. That the behaviour is unchanged is the whole point --- the application
cannot tell an FPGA from a stub.

\solution{ex:12:4}{Freestanding}
The exercise asks you to build; the questions are answered here.
\ans{a} Yes. Nothing in the project uses exceptions or RTTI.
\ans{b} Nothing should be found. The whole stack --- \code{Frame}, \code{Interface}, \code{Stub},
\code{ByteTransport}, \code{Spi}, \code{EchoNode} --- is built on statically allocated objects of
known size.
\ans{c} Because the coding conventions were written for a freestanding target from the start:
return codes instead of exceptions, aggregates instead of dynamic containers, virtual functions
instead of downcasting. That was not luck, and this is the chapter that shows what the habit was
worth.

\solution{ex:12:5}{The construction order}
\ans{a} It compiles, and the compiler does not warn. Members are constructed in
\emph{declaration order}, whatever order the initialiser list happens to have.
\ans{b} \code{myDriver} is constructed first and binds a reference to \code{myTransport}, which has
not been constructed yet. The binding itself is harmless --- it is only an address --- but the
first call to \code{transfer()} uses an object that has not been initialised, that is, undefined
behaviour. On a freestanding target it usually manifests as the SPI peripheral never being
configured.
\ans{c} That a member depending on another has to be declared \emph{after} it. The initialiser
list does not express the order; it just follows the declaration order, and a compiler that warns
about it only does so when the list and the declarations are in different orders --- not here.
